\documentclass[10pt]{article}
\usepackage{logo}
\usepackage{fullpage}
\usepackage[utf8]{inputenc}
\usepackage[spanish]{babel}
\usepackage{epsfig}
\usepackage{amsmath}
\usepackage{amssymb}
\usepackage{multicol}
\usepackage{color}

\newcommand{\cyan}[1]{\textCyan #1\textBlack}
\newcommand{\red}[1]{\textRed #1\textBlack}
\newcommand{\green}[1]{\textGreen #1\textBlack}
\newcommand{\blue}[1]{\textBlue #1\textBlack}
\newcommand{\black}[1]{\textBlack #1\textBlack}
\newcommand{\magenta}[1]{\textMagenta #1\textBlack}
\newcommand{\brown}[1]{\textBrown #1\textBlack}
\newcommand{\vs}[1]{\vspace{#1mm}}
\newcommand{\ignore}{}
\newcommand{\ri}[1]{\text{\red{#1}}}
\newcommand{\hs}{\hat\sigma}
\newcommand{\modelos}{{\it modelos}}
\newcommand{\B}{{\tt B}}
\newcommand{\nspace}{\text{NSPACE}}
\newcommand{\logspace}{\text{LOGSPACE}}
\newcommand{\nlogspace}{\text{NLOGSPACE}}
\newcommand{\npspace}{\text{NPSPACE}}
\newcommand{\pspace}{\text{PSPACE}}
\newcommand{\ph}{\text{PH}}
\newcommand{\expspace}{\text{EXPSPACE}}
\newcommand{\nexpspace}{\text{NEXPSPACE}}
\newcommand{\dspace}{\text{DSPACE}}
\newcommand{\espacio}{{\it espacio}}
\newcommand{\tiempo}{{\it tiempo}}
\newcommand{\ptime}{\text{PTIME}}
\newcommand{\dtime}{\text{DTIME}}
\newcommand{\exptime}{\text{EXPTIME}}
\newcommand{\nexptime}{\text{NEXPTIME}}
\newcommand{\CC}{{\cal C}}
\newcommand{\sat}{\text{SAT}}
\newcommand{\tcnfu}{\text{3-CNF-SAT-UNSAT}}
\newcommand{\usat}{\text{unique-SAT}}
\newcommand{\np}{\text{NP}}
\newcommand{\ntime}{\text{NTIME}}
\newcommand{\cnf}{\text{CNF-SAT}}
\newcommand{\tcnf}{\text{3-CNF-SAT}}
\newcommand{\dcnf}{\text{2-CNF-SAT}}
\newcommand{\horn}{\text{HORN-SAT}}
\newcommand{\nhorn}{\text{NEG-HORN-SAT}}
\newcommand{\co}{\text{co-}}
\newcommand{\rp}{\leq^\text{\it p}_\text{\it m}}
\newcommand{\tur}{\leq^\text{\it p}_\text{\it T}}
\newcommand{\reach}{\text{CAMINO}}
\newcommand{\pe}{\text{PROG-ENT}}
\newcommand{\pl}{\text{PROG-LIN}}
\newcommand{\cor}{\text{CONT-REG}}
\newcommand{\er}{\text{EQUIV-REG}}
\newcommand{\qbf}{\text{QBF}}
\newcommand{\shp}{\text{Succinct-HP}}
\newcommand{\hp}{\text{HP}}
\newcommand{\cdp}{\text{DP}}
\newcommand{\clique}{\text{CLIQUE}}
\newcommand{\eclique}{\text{exact-CLIQUE}}
\newcommand{\costo}{\text{costo}}
\newcommand{\tsp}{\text{TSP}}
\newcommand{\tspu}{\text{unique-TSP}}
\newcommand{\no}{\text{NO}}
\newcommand{\yes}{\text{YES}}
\newcommand{\br}{\text{CERTAIN-ANSWERS}}
\newcommand{\dpc}{\text{DP}}

\newcommand{\CROM}{\text{CROM}}
\newcommand{\EVAL}{\text{EVAL}}
\newcommand{\EQUIV}{\text{EQUIV}}

\newcommand{\re}{\text{RE}}
\newcommand{\dec}{\text{R}}
\newcommand{\resp}[1]{{\text{\bf [Responsable: #1]}}}
\usepackage{braket}
\usepackage{pythonhighlight}
% pythonhighlight activa \UseRawInputEncoding; restauramos UTF-8 para el resto del documento.
\inputencoding{utf8}
% babel-spanish redefine \% como comando robusto, lo que produce el error
% "Improper alphabetic constant" en el entorno python de pythonhighlight.
\addto\extrasspanish{\ifdefined\percentsign\let\%\percentsign\fi}
\usepackage{url}
\usepackage[colorlinks=true,urlcolor=blue,linkcolor=black]{hyperref}

\begin{document}


\begin{tabular}{ccl}
\begin{tabular}{c}
\psfig{file=puclogo.eps}
\end{tabular}
&\ \ \ &
\begin{tabular}{l}
PONTIFICIA UNIVERSIDAD CATÓLICA DE CHILE\\
ESCUELA DE INGENIERÍA\\
DEPARTAMENTO DE CIENCIA DE LA COMPUTACIÓN
\end{tabular}
\end{tabular}

\vspace{1cm}

\begin{center}
  {
\bf Tópicos Avanzados en Teoría de la Computación - IIC3810\\
\bf Tarea 3\\
\bf Fecha de entrega: jueves 19 de noviembre}
\end{center}

\vspace{5mm}

\section*{Preguntas}

\begin{enumerate}
\item \textbf{[2 puntos]} En esta pregunta usted deberá implementar en Python una librería
  que permita utilizar el algoritmo de estimación de fase, o QPE. Esta librería debe estar
  en el archivo \texttt{qpe.py}, y se utilizará tanto en un simulador como en un computador
  cuántico de IBM. Las funciones de esta librería reciben una matriz unitaria
  $U\in\mathbb C^{2^n\times2^n}$, un vector unitario $\ket{w}\in\mathbb C^{2^n}$
  que es vector propio de $U$, y un entero $t\geq1$ que indica el número de qubits que se
  utilizarán para estimar la fase. Los parámetros \texttt{U} y \texttt{w} deben ser
  arreglos de números complejos de NumPy, de tipo \texttt{np.complex128}. Puede suponer que $n\geq1$ y que las entradas
  satisfacen las condiciones anteriores. Su librería debe implementar las siguientes funciones:

  \begin{itemize}
  \item Una función \texttt{circuito\_qpe} con el siguiente encabezado:
\begin{python}
def circuito_qpe(U: np.ndarray, w: np.ndarray, t: int) -> QuantumCircuit:
\end{python}
    Esta función retorna un objeto de la clase \texttt{QuantumCircuit} que implementa
    QPE para la entrada indicada, con $n+t$ qubits y $t$ bits clásicos. El circuito
    comienza con todos los qubits en cero y debe incluir la preparación del estado
    $\ket{0^t}\otimes\ket{w}$. Al final debe medir únicamente el registro de fase,
    almacenando el resultado en los $t$ bits clásicos, de modo que el circuito pueda
    ejecutarse tanto en el simulador como en el computador cuántico.

  \item Una función \texttt{ejecutar\_qpe} con el siguiente encabezado:
\begin{python}
def ejecutar_qpe(U: np.ndarray, w: np.ndarray, t: int,
                                num_iter: int = 10000) -> dict[int, int]:
\end{python}
    Esta función ejecuta \texttt{num\_iter} veces el circuito \texttt{circuito\_qpe(U, w, t)}
    en un simulador ideal, utilizando \texttt{StatevectorSampler} de \texttt{qiskit.primitives}, y retorna
    un diccionario cuyas claves son los enteros $\ell\in\{0,\ldots,2^t-1\}$ observados
    y cuyos valores indican el número de veces que se obtuvo cada resultado. La estimación
    asociada a la clave $\ell$ es $\ell/2^t$. Puede suponer que \texttt{num\_iter} es un
    entero positivo, cuyo valor por defecto es 10000, y la suma de los valores del
    diccionario debe ser igual a este parámetro.
  \end{itemize}

\item \textbf{[2 puntos]} En esta pregunta usted deberá utilizar su librería para estimar
  una fase irracional en un simulador de Qiskit y estudiar experimentalmente la garantía
  de error de QPE. Para esto, considere la matriz $U\in\mathbb C^{8\times8}$ que se
  define a continuación a partir de los números complejos
  \begin{align*}
    a&=\frac{-4-\sqrt{2+\sqrt2}+2\cos\!\left(\frac{\pi}{\sqrt2}\right)
      +i\left[\sqrt{2-\sqrt2}+2\sin\!\left(\frac{\pi}{\sqrt2}\right)\right]}{16},\\[3pt]
    b&=\frac{-\sqrt{2+\sqrt2}-2\cos\!\left(\frac{\pi}{\sqrt2}\right)
      +i\left[\sqrt{2-\sqrt2}-2\sin\!\left(\frac{\pi}{\sqrt2}\right)\right]}{16},\\[3pt]
    d&=\frac{12-\sqrt{2+\sqrt2}+2\cos\!\left(\frac{\pi}{\sqrt2}\right)
      +i\left[\sqrt{2-\sqrt2}+2\sin\!\left(\frac{\pi}{\sqrt2}\right)\right]}{16}.
  \end{align*}
  La matriz $U$ y el vector $\ket{w}$ que utilizará en esta pregunta son
  \[
  \begingroup
  \setlength{\arraycolsep}{4pt}
  \renewcommand{\arraystretch}{0.95}
  U=\begin{pmatrix}
    d&b&a&b&a&b&a&b\\
    b&d&b&a&b&a&b&a\\
    a&b&d&b&a&b&a&b\\
    b&a&b&d&b&a&b&a\\
    a&b&a&b&d&b&a&b\\
    b&a&b&a&b&d&b&a\\
    a&b&a&b&a&b&d&b\\
    b&a&b&a&b&a&b&d
  \end{pmatrix},
  \qquad
  \ket{w}=\frac1{\sqrt8}
  \begin{pmatrix}
    1\\-1\\1\\-1\\1\\-1\\1\\-1
  \end{pmatrix}.
  \endgroup
  \]
La matriz $U$ es unitaria y el vector $\ket{w}$ satisface
  $U\ket{w}=e^{2\pi i\theta}\ket{w}$, donde $\theta=\sqrt2/4$.
  Usted debe escribir un programa \texttt{pregunta2.py} que construya \texttt{U} y
  \texttt{w} evaluando estas expresiones con \texttt{np.sqrt}, \texttt{np.cos},
  \texttt{np.sin}, \texttt{np.pi} y la unidad imaginaria \texttt{1j}. Al crear los
  arreglos debe indicar \texttt{dtype=np.complex128}; sus entradas serán aproximaciones
  numéricas de las expresiones dadas. Las variables \texttt{U} y \texttt{w}
  se deben definir en el nivel superior del archivo, y la ejecución del experimento
  debe estar dentro de un bloque
  \mbox{\texttt{if \_\_name\_\_ == "\_\_main\_\_":}}. De esta forma, la matriz y el
  vector podrán ser importados desde otro archivo sin ejecutar el experimento, lo que
  será necesario en la pregunta 3.

  Su programa \texttt{pregunta2.py} debe importar \texttt{ejecutar\_qpe} desde
  \texttt{qpe.py} y realizar el experimento para cada una de las cuatro
  combinaciones de $m\in\{3,7\}$ y
  $\varepsilon\in\{1/8,1/32\}$. Para cada combinación, debe determinar $t$
  a partir de la cota de error de QPE, para una tolerancia $2^{-m}$ y una probabilidad
  de error a lo más $\varepsilon$, y calcular $Q=2^t$. Luego debe utilizar
  \texttt{ejecutar\_qpe} con \texttt{U}, \texttt{w}, el valor de $t$ obtenido y
  $\texttt{num\_iter}=10000$.
  A partir de este diccionario, el programa debe calcular la fracción $p_{\mathrm{sim}}$ de los resultados $\ell$ cuya estimación $\ell/Q$ está a distancia a lo más $2^{-m}$ de $\theta$. Además,
  debe entregar una única estimación $\ell^\star/Q$ de la fase, para lo cual usted debe decidir qué valor $\ell$ utilizar entre todos los que aparecieron al menos una vez.

  Para esta pregunta, debe entregar únicamente el programa \texttt{pregunta2.py}.
  Al ejecutarse, el programa debe imprimir exactamente cuatro líneas, en el orden y
  con el formato que se muestra a continuación, reemplazando los puntos suspensivos
  por los resultados calculados. En cada línea, la frecuencia de éxito corresponde
  a $p_{\mathrm{sim}}$ y debe escribirse con cuatro cifras decimales; la fase estimada
  debe escribirse como la fracción $\ell^\star/Q$, con numerador y denominador enteros.

\begin{verbatim}
Para m = 3 y epsilon = 1/8: frecuencia de éxito = ...; fase estimada = .../...
Para m = 3 y epsilon = 1/32: frecuencia de éxito = ...; fase estimada = .../...
Para m = 7 y epsilon = 1/8: frecuencia de éxito = ...; fase estimada = .../...
Para m = 7 y epsilon = 1/32: frecuencia de éxito = ...; fase estimada = .../...
\end{verbatim}

\item \textbf{[2 puntos]} En la Tarea 2 usted ejecutó el algoritmo de Deutsch-Jozsa en
  un computador cuántico de IBM. En esta pregunta deberá ejecutar QPE en uno de estos
  computadores. Para esto, debe utilizar la misma matriz $U$, el mismo
  vector $\ket{w}$ y las mismas cuatro combinaciones de $m$ y $\varepsilon$ de la
  pregunta anterior.
En particular, usted debe escribir un programa \texttt{pregunta3.py} que importe
  \texttt{U} y \texttt{w} desde \texttt{pregunta2.py}, y la función
  \texttt{circuito\_qpe} desde \texttt{qpe.py}.
  Para cada combinación de $m$ y $\varepsilon$, el programa debe utilizar el valor
  de $t$ obtenido en la pregunta 2 para construir \texttt{circuito\_qpe(U, w, t)},
  transpilarlo y ejecutarlo con 1000 shots. Las cuatro combinaciones producen cuatro
  circuitos distintos.

  Un trabajo (\emph{job}) es una solicitud que se envía a IBM para ejecutar uno o más
  circuitos. En esta pregunta, usted debe enviar un trabajo separado por circuito,
  para poder conocer el tiempo de uso de QPU correspondiente a cada uno. Para esto,
  puede utilizar \texttt{sampler.run([circuito], shots=1000)}, donde \texttt{circuito}
  es uno de los circuitos transpilados. Esta llamada solicita 1000 ejecuciones de ese
  circuito, por lo que debe enviar cuatro trabajos y realizará 4000 ejecuciones en total.
  Para cada circuito, guarde un diccionario en el que la clave $\ell$ sea un resultado
  observado y el valor $c_\ell$ sea el número de veces que apareció en esas 1000
  ejecuciones; por lo tanto, la suma de los valores del diccionario debe ser 1000.
  Este diccionario contiene los conteos que utilizará para analizar ese circuito.

  Para cada combinación de $m$ y $\varepsilon$, a partir del diccionario correspondiente,
  el programa debe calcular la fracción $p_{\mathrm{IBM}}$ de los resultados $\ell$
  cuya estimación $\ell/Q$ está a distancia a lo más $2^{-m}$ de $\theta$. Además,
  debe entregar una única estimación $\ell^\star/Q$ de la fase, para lo cual usted
  debe decidir qué valor $\ell$ utilizar entre todos los que aparecieron al menos una vez.

  Al ejecutarse, el programa debe imprimir exactamente cuatro líneas, en el orden y
  con el formato que se muestra a continuación, reemplazando los puntos suspensivos
  por los resultados calculados. En cada línea, la frecuencia de éxito corresponde
  a $p_{\mathrm{IBM}}$ y debe escribirse con cuatro cifras decimales; la fase estimada
  debe escribirse como la fracción $\ell^\star/Q$, con numerador y denominador enteros.

\begin{verbatim}
Para m = 3 y epsilon = 1/8: frecuencia de éxito = ...; fase estimada = .../...
Para m = 3 y epsilon = 1/32: frecuencia de éxito = ...; fase estimada = .../...
Para m = 7 y epsilon = 1/8: frecuencia de éxito = ...; fase estimada = .../...
Para m = 7 y epsilon = 1/32: frecuencia de éxito = ...; fase estimada = .../...
\end{verbatim}

  Para esta pregunta, debe entregar el programa \texttt{pregunta3.py} y el archivo
  \texttt{resultados3.pdf}. Este archivo debe contener las cuatro líneas de salida
  del programa, con las frecuencias de éxito y las fases estimadas obtenidas al
  ejecutar los circuitos en el computador cuántico de IBM. Además, debe indicar el
  tiempo de uso de QPU reportado por IBM para cada uno de los cuatro trabajos, en
  segundos, identificando la combinación de $m$ y $\varepsilon$ correspondiente. Puede consultar este tiempo con \texttt{job.usage()}
  una vez que el trabajo haya finalizado. Este valor corresponde al uso de QPU del
  trabajo completo, incluidas las 1000 ejecuciones del circuito, y excluyendo la espera en la cola de trabajo.
\end{enumerate}

\section*{Nota sobre la entrega y corrección}

La entrega de esta tarea debe incluir los tres programas \texttt{qpe.py},
\texttt{pregunta2.py} y \texttt{pregunta3.py}, junto con el documento
\texttt{resultados3.pdf}, que debe contener la salida de \texttt{pregunta3.py} y
los tiempos de uso de QPU solicitados en la pregunta 3.
La entrega de la pregunta 2 consiste únicamente en
el programa \texttt{pregunta2.py}. Verifique los circuitos en un
simulador antes de enviarlos a IBM y conserve los conteos obtenidos para consultarlos
sin repetir las ejecuciones.

En la corrección se revisará su código y la salida incluida en \texttt{resultados3.pdf};
\texttt{pregunta3.py} no se ejecutará en un computador cuántico. Debe reportar todos
los resultados obtenidos, incluso si alguna frecuencia de éxito es menor que
$1-\varepsilon$. No se exigirá alcanzar esta
cota en el hardware; la evaluación considerará la implementación del experimento
y los resultados reportados.

\end{document}
